\documentclass[11pt]{article}

\usepackage[margin=0.95in]{geometry}
\usepackage{amsmath,amssymb,amsthm,mathtools,booktabs,enumitem,microtype}
\usepackage{xcolor}
\usepackage[colorlinks=true,linkcolor=blue!50!black,urlcolor=blue!50!black,citecolor=blue!50!black]{hyperref}

\theoremstyle{plain}
\newtheorem{theorem}{Theorem}
\newtheorem{lemma}[theorem]{Lemma}
\newtheorem{proposition}[theorem]{Proposition}
\newtheorem{corollary}[theorem]{Corollary}

\theoremstyle{definition}
\newtheorem{definition}[theorem]{Definition}

\theoremstyle{remark}
\newtheorem{remark}{Remark}

\newcommand{\RL}{\mathrm{RL}}
\newcommand{\Int}{\mathrm{Int}}

\title{\bfseries Optimality of the Greedy Cutting Schedule for the Randomized Littlestone Dimension of Intervals}
\author{}
\date{\vspace{-1.6em}%
\small
\begin{tabular}{@{}c@{}}
Target paper: M.~Daga and S.~P.~Ramu, \emph{How many mistakes does it take to learn an interval?}\\
\emph{Exact Littlestone dimensions of intervals and their unions},\\
\href{https://arxiv.org/abs/2611.28727}{arXiv:2611.28727} (November 2026), Section~7, Conjecture~\texttt{conj:rl}\\[0.35em]
Companion verification script: \texttt{solutions/interval\_littlestone/explore.py}\\[0.35em]
November 2026
\end{tabular}}

\begin{document}

\maketitle

\begin{abstract}
We prove Conjecture~\texttt{conj:rl} of Daga and Ramu (\href{https://arxiv.org/abs/2611.28727}{arXiv:2611.28727}, November 2026): for every finite multiset $S$ of positive block lengths, the greedy canonical cutting schedule attains the exact randomized Littlestone dimension $\RL(J_S)$, and consequently $\RL(\Int_m) = R(m + 1)$ for every integer $m \ge 0$.
The proof decomposes the piecewise-linear logarithm $\tilde\lambda(x)$ into $[0,1]$-valued dyadic layers $h_d(x) = Q_d(x) - 2Q_{2d}(x)$ ($d = 2^k$).
At every dyadic scale, each single-rod canonical tree is \emph{packed} with total mass $Q_d(x)$, and packed majorization bounds the continuation gain from any non-canonical split in the presence of arbitrary unaffected rods by $\frac{1}{2}\sum_d E_d$, which telescopes identically to the immediate-reward loss.
A closed-form sum of packed layer values then directly recovers the single-interval recurrence $R(p)$.
\end{abstract}

\section{Reduction to Discounted Rod Cutting}

Given a finite multiset $S$ of positive integers, place one disjoint ordered block of $g$ points for each occurrence of $g \in S$. The class $J_S$ consists of the empty set and all nonempty intervals contained in a single block; in particular $\Int_m = J_{\{m\}}$ (containing only the empty hypothesis when $m = 0$). Represent each block of $g$ points by a \emph{rod} of measure $p := g + 1 \ge 2$; unit rods ($p = 1$) carry no queryable points and may be retained or omitted freely. All unions of rods below are multiset unions.

For a nonempty finite hypothesis class $H$, its randomized Littlestone dimension $\RL(H)$ is half the maximum expected branch length of a binary query tree shattered by $H$ under uniform child choices, obeying
\[
\RL(H) = \max\Bigl(\{0\} \cup \Bigl\{\tfrac{1}{2}\bigl(1 + \RL(H_x^0) + \RL(H_x^1)\bigr) : x \text{ splits } H\Bigr\}\Bigr),
\]
where $H_x^b := \{h \in H : h(x) = b\}$ must both be nonempty.
On disjoint domains, $\RL(H_1 \otimes H_2) = \RL(H_1) + \RL(H_2)$ by induction on $|H_1| + |H_2|$ (every query splits one factor while leaving the other intact).

Following Section~7 of~\cite{DagaRamu2026}, denote by $\tilde\lambda$ the piecewise-linear interpolation of $\log_2$ at powers of two:
\[
\tilde\lambda(x) := t + \frac{x - 2^t}{2^t} \qquad (x \ge 1,\; t = \lfloor \log_2 x \rfloor).
\]
A chain of $c$ endpoint thresholds satisfies $\RL = \tfrac{1}{2}\tilde\lambda(c)$: every split leaves chains of sizes $u, v \ge 1$ with $u + v = c$; concavity and $\tilde\lambda(2x) = 1 + \tilde\lambda(x)$ give $\tilde\lambda(u) + \tilde\lambda(v) \le 2\tilde\lambda(c) - 2$, attained by $\lfloor c/2 \rfloor$ and $\lceil c/2 \rceil$ since both lie in one affine piece of $\tilde\lambda$.

Querying the $u$-th point of a block of rod measure $p$ ($1 \le u \le p - 1$, $v := p - u \ge 1$) produces two independent endpoint chains of sizes $u$ and $v$ on the positive branch and splits the rod into $\{u, v\}$ on the negative branch while leaving all other rods $T$ unchanged. Writing $V(P) := \RL(J_S)$ for the corresponding rod multiset $P$, the Bellman recursion (\texttt{eq:rlcut} in~\cite{DagaRamu2026}) reads
\begin{equation}\label{eq:bellman}
V(P) = \max_{\substack{p \in P,\, p \ge 2 \\ u + v = p,\; u, v \ge 1}}
\frac{1}{2}\Bigl(c(u, v) + V\bigl((P \setminus \{p\}) \cup \{u, v\}\bigr)\Bigr),
\qquad
c(u, v) := 1 + \frac{\tilde\lambda(u) + \tilde\lambda(v)}{2},
\end{equation}
with $V(\varnothing) = 0$. Every legal cut decreases the rank $\sum_{r \in P}(r - 1)$ by~$1$.

For $p \ge 2$, set $b := 2^{\lfloor \log_2 p \rfloor - 1}$ and define the \emph{canonical split} of $p$ by
\[
a := \max(b,\, p - 2b), \qquad a' := p - a.
\]
Both $a$ and $a'$ lie in $[b, 2b]$, and at least one equals $b$ or $2b$. Affineness of $\tilde\lambda$ on $[b, 2b]$ together with $\tilde\lambda(2x) = 1 + \tilde\lambda(x)$ gives $c(a, a') = \tilde\lambda(p)$.
The \emph{greedy schedule} always applies the canonical split to a largest available rod. Recursively splitting canonically down to units defines a finite binary \emph{canonical forest} for any multiset $P$. Since children are strictly smaller than parents, greedy visits all internal nodes of this forest in non-increasing order of measure (breaking ties arbitrarily). Hence its discounted value $G(P)$ is $\sum_{i \ge 1} 2^{-i} \tilde\lambda(p_i^\downarrow)$, where $p_1^\downarrow \ge p_2^\downarrow \ge \cdots$ lists the internal-node measures of the canonical forest of $P$.

\section{Dyadic Layer Decomposition and Packedness}

For each dyadic scale $d \in \{1, 2, 4, \dots\}$ and $x \ge 1$, define
\[
Q_d(x) := \max\!\left(0,\, \frac{x}{d} - 1\right), \qquad
h_d(x) := Q_d(x) - 2Q_{2d}(x) = \min\bigl(1,\, Q_d(x)\bigr).
\]
The second equality is immediate on each of the three regimes $x \le d$, $d < x \le 2d$, and $x > 2d$. When $2^t \le x < 2^{t+1}$, we have $h_d(x) = 1$ for $d = 1, \dots, 2^{t-1}$, $h_{2^t}(x) = (x - 2^t)/2^t$, and $h_d(x) = 0$ for $d > 2^t$. Summing both expressions for $h_d(x)$ over all dyadic $d$ (finite support at every fixed $x$) yields
\begin{equation}\label{eq:dyadic-sums}
\tilde\lambda(x) = \sum_d h_d(x), \qquad
\sum_d Q_d(x) = 2(x - 1) - \tilde\lambda(x).
\end{equation}

For any finite list $X$ of nonnegative numbers, define $D(X) := \sum_{i \ge 1} 2^{-i} X_i^\downarrow$ with entries sorted descending and zero-padded. Let $A_d(P)$ denote the multiset of values $h_d(r)$ over all internal nodes $r$ of the canonical forest of $P$. Because each $h_d$ is nondecreasing in $r$, one common non-increasing node ordering sorts every layer simultaneously, whence
\begin{equation}\label{eq:layer-decomp}
G(P) = \sum_d D\bigl(A_d(P)\bigr).
\end{equation}

Call a list with entries in $[0, 1]$ \emph{packed} if, after sorting in descending order and dropping zeros, it consists of ones followed by at most one entry in $(0, 1)$.

\begin{lemma}[Single-rod packedness]\label{lem:single-packed}
For every integer $x \ge 1$ and every dyadic $d = 2^k$ ($k \ge 0$), the list $A_d(\{x\})$ is packed and has total mass (sum of entries) $Q_d(x)$.
\end{lemma}

\begin{proof}
Induct on $x$. If $x \le d$, every node in the canonical tree of $x$ has measure $\le d$, so $A_d(\{x\})$ vanishes and has mass $0 = Q_d(x)$. If $d < x < 2d$ (which requires $d \ge 2$), then $b = d/2$, so both canonical children lie in $[d/2, d]$ and contribute zero; only the root contributes, with entry $h_d(x) = x/d - 1 \in (0, 1)$, which is packed of mass $Q_d(x)$.

If $x \ge 2d$, then $b = 2^{\lfloor \log_2 x \rfloor - 1} \ge d$, so both canonical children $a, a'$ satisfy $a, a' \ge b \ge d$, and at least one is a power of two $\ge b \ge d$, hence divisible by $d$. The root contributes $h_d(x) = 1$, and
\[
1 + Q_d(a) + Q_d(a') = 1 + \left(\frac{a}{d} - 1\right) + \left(\frac{a'}{d} - 1\right) = \frac{x}{d} - 1 = Q_d(x).
\]
By induction, $A_d(\{a\})$ and $A_d(\{a'\})$ are packed with masses $Q_d(a)$ and $Q_d(a')$. The child divisible by $d$ has integer mass, so its positive entries are all ones. Combining both child lists with the root entry~$1$ produces a packed list of mass $Q_d(x)$.
\end{proof}

\begin{lemma}[Per-layer split comparison]\label{lem:split-layer}
Let $p \ge 2$, let $a + a' = p$ be its canonical split, and let $u + v = p$ ($u, v \ge 1$) be any legal split. For every dyadic $d = 2^k$, the list $A_d(\{a, a'\})$ is packed with mass $M_d := Q_d(a) + Q_d(a')$, and $A_d(\{u, v\})$ has mass $M_d + E_d$ for some excess $E_d \ge 0$.
\end{lemma}

\begin{proof}
Since $p \ge 2$, $b = 2^{\lfloor \log_2 p \rfloor - 1} \ge 1$ is a power of two, so either $d \le b$ or $d \ge 2b$.
If $d \le b$, then $a, a' \ge b \ge d$ and one canonical child is a power of two divisible by $d$; by Lemma~\ref{lem:single-packed}, its layer has only unit positive entries, so $A_d(\{a, a'\})$ is packed with mass
\[
M_d = \left(\frac{a}{d} - 1\right) + \left(\frac{a'}{d} - 1\right) = \frac{p}{d} - 2 \le Q_d(u) + Q_d(v) = M_d + E_d.
\]
If $d \ge 2b$, then $a, a' \le 2b \le d$, so $A_d(\{a, a'\})$ vanishes ($M_d = 0$, trivially packed) and $E_d = Q_d(u) + Q_d(v) \ge 0$.
\end{proof}

\section{Packed Majorization and Optimality of Greedy}

\begin{lemma}[Packed majorization under union and discounting]\label{lem:majorize}
Let $A$ be a packed list of mass $M$, let $B$ be any list with entries in $[0, 1]$ of mass $M + E$ ($E \ge 0$), and let $C$ be any list with entries in $[0, 1]$. Then
\begin{equation}\label{eq:majorize}
D(C \cup B) - D(C \cup A) \;\le\; \frac{E}{2}.
\end{equation}
\end{lemma}

\begin{proof}
For any nonnegative list $X$ and integer $k \ge 0$, let $s_k(X) := \sum_{i=1}^k X_i^\downarrow$ (with zero padding and $s_0 = 0$). Because $A$ is packed of mass $M$, $s_k(A) = \min(k, M)$ for all $k \ge 0$. Because $B$ has entries in $[0, 1]$ summing to $M + E$,
\[
s_k(B) \le \min(k, M + E) \le \min(k, M) + E = s_k(A) + E \qquad (\forall k \ge 0).
\]
Using the exact prefix-sum identity
\[
s_k(C \cup X) = \max_{0 \le j \le k} \bigl\{s_j(C) + s_{k - j}(X)\bigr\},
\]
we obtain $s_k(C \cup B) \le s_k(C \cup A) + E$ for every $k \ge 1$.
By Abel summation, $D(X) = \sum_{k \ge 1} 2^{-k-1} s_k(X)$ (each entry $X_i^\downarrow$ receives weight $\sum_{k \ge i} 2^{-k-1} = 2^{-i}$). Since $\sum_{k \ge 1} 2^{-k-1} = 1/2$,
\[
D(C \cup B) - D(C \cup A) = \sum_{k \ge 1} 2^{-k-1}\bigl(s_k(C \cup B) - s_k(C \cup A)\bigr) \le \frac{E}{2}. \qedhere
\]
\end{proof}

\begin{theorem}[Optimality of the greedy cutting schedule]\label{thm:greedy}
For every finite multiset $P$ of rods, $V(P) = G(P)$. In particular, the greedy schedule attains $\RL(J_S)$ for every finite multiset $S$ of positive block lengths.
\end{theorem}

\begin{proof}
Induct on the rank $\sum_{r \in P}(r - 1)$. At rank~$0$, $V(P) = G(P) = 0$.
In positive rank, consider any legal first cut splitting $p \in P$ ($p \ge 2$) into $u + v = p$ ($u, v \ge 1$), write $T := P \setminus \{p\}$, and let $a + a' = p$ be the canonical split of $p$.
Applying Lemma~\ref{lem:majorize} at every dyadic scale $d$ with $C = A_d(T)$, $A = A_d(\{a, a'\})$, and $B = A_d(\{u, v\})$, and summing via~\eqref{eq:layer-decomp} and Lemma~\ref{lem:split-layer}, yields
\[
G(T \cup \{u, v\}) - G(T \cup \{a, a'\}) \;\le\; \frac{1}{2}\sum_d E_d.
\]
By~\eqref{eq:dyadic-sums} and $u + v = a + a' = p$, the right-hand side telescopes identically to the immediate-reward loss:
\begin{align}
\frac{1}{2}\sum_d E_d
&= \frac{1}{2}\sum_d \bigl(Q_d(u) + Q_d(v) - Q_d(a) - Q_d(a')\bigr) \nonumber \\
&= \frac{1}{2}\bigl(\tilde\lambda(a) + \tilde\lambda(a') - \tilde\lambda(u) - \tilde\lambda(v)\bigr)
= c(a, a') - c(u, v). \label{eq:cancel}
\end{align}
By the induction hypothesis, $V(T \cup \{u, v\}) = G(T \cup \{u, v\})$. Combining~\eqref{eq:bellman}, \eqref{eq:cancel}, and $c(a, a') = \tilde\lambda(p)$, the cut $(u, v)$ has value at most
\[
\frac{1}{2}\bigl(c(u, v) + G(T \cup \{u, v\})\bigr)
\;\le\;
\frac{1}{2}\bigl(\tilde\lambda(p) + G(T \cup \{a, a'\})\bigr).
\]
The right-hand side is the discounted value of cutting $p$ first and then processing all remaining internal nodes of the canonical forest of $T \cup \{p\}$ in non-increasing order of measure.
Since the discount weights $2^{-i}$ are decreasing and $\tilde\lambda$ is nondecreasing, sorting \emph{all} internal nodes of the canonical forest of $T \cup \{p\}$ in non-increasing order gives at least as much, and this sorted order is feasible and equals greedy on $P = T \cup \{p\}$.
Hence $V(P) \le G(P)$, and since greedy is feasible, $V(P) = G(P)$.
\end{proof}

\section{Exact Recurrence for a Single Interval Class}

\begin{theorem}[Single-interval closed form]\label{thm:single}
For every integer $p \ge 1$, $G(\{p\}) = R(p)$, where $R(1) = 0$ and
\[
R(p) = 1 - 2^{1 - p} + \frac{1}{2}\bigl(R(\lfloor p/2 \rfloor) + R(\lceil p/2 \rceil)\bigr) \qquad (p \ge 2).
\]
Consequently, $\RL(\Int_m) = R(m + 1)$ for every integer $m \ge 0$.
\end{theorem}

\begin{proof}
Define $H(y) := 0$ for $0 \le y \le 1$, set $H(n) := 1 - 2^{1 - n}$ for every integer $n \ge 1$, and interpolate $H$ linearly between consecutive integers.
For $y = n + \theta$ with $n = \lfloor y \rfloor \ge 1$ and $\theta \in [0, 1)$, a packed list of mass $y - 1$ consists of $n - 1$ ones followed by $\theta$, whose discounted sum is
\[
\sum_{i = 1}^{n - 1} 2^{-i} + \theta\,2^{-n} = (1 - 2^{1 - n}) + \theta\,2^{-n} = H(y).
\]
By Lemma~\ref{lem:single-packed} and~\eqref{eq:layer-decomp},
\begin{equation}\label{eq:G-sum}
G(\{p\}) = \sum_{j = 0}^\infty H(p / 2^j) \qquad (\forall p \ge 1).
\end{equation}
For $p = 1$, every term vanishes, so $G(\{1\}) = 0 = R(1)$. For $p \ge 2$, separating the $j = 0$ term $H(p) = 1 - 2^{1 - p}$ from~\eqref{eq:G-sum} gives
\[
G(\{p\}) = 1 - 2^{1 - p} + \sum_{j = 0}^\infty H\!\left(\frac{p/2}{2^j}\right).
\]
For every $j \ge 0$, the function $x \mapsto H(x / 2^j)$ has breakpoints only at multiples of $2^j \in \mathbb{Z}_{\ge 1}$, hence is affine on $[\lfloor p/2 \rfloor, \lceil p/2 \rceil]$, whose midpoint is $p/2$. Thus
\[
H\!\left(\frac{p/2}{2^j}\right) = \frac{1}{2}\left(H\!\left(\frac{\lfloor p/2 \rfloor}{2^j}\right) + H\!\left(\frac{\lceil p/2 \rceil}{2^j}\right)\right) \qquad (\forall j \ge 0).
\]
Summing over $j \ge 0$ shows that $G(\{p\})$ satisfies the defining recurrence of $R(p)$. By Theorem~\ref{thm:greedy}, $\RL(\Int_m) = V(\{m + 1\}) = G(\{m + 1\}) = R(m + 1)$ for all $m \ge 0$.
\end{proof}

\begin{remark}[Lean~4 formalization compatibility]
Daga and Ramu~\cite{DagaRamu2026} machine-check their deterministic dimension formulas in Lean~4. Because Lemmas~\ref{lem:single-packed}--\ref{lem:majorize} and Theorems~\ref{thm:greedy}--\ref{thm:single} use only finitary dyadic arithmetic and finite sums of rational numbers, the proof integrates naturally into the same Lean~4 formalization.
\end{remark}

\section{Exact Computational Verification}

The script \texttt{solutions/interval\_littlestone/explore.py} verifies Theorems~\ref{thm:greedy} and~\ref{thm:single} in exact integer and rational arithmetic.
It confirms that the exact Bellman optimum of~\eqref{eq:bellman} equals the greedy heap schedule on all $204{,}226$ multisets of non-unit rods with total measure at most~$50$ (evaluating $2{,}977{,}224$ distinct legal first cuts), matches $R(m + 1)$ on single rods through $m = 49$ (unrestricted Bellman DP) and $m = 511$ (direct greedy simulation), and agrees with the explicit shattered-tree recursion on $\Int_m$ for all $0 \le m \le 10$. All assertions pass.

\begin{thebibliography}{9}

\bibitem{DagaRamu2026}
M.~Daga and S.~P.~Ramu.
\emph{How many mistakes does it take to learn an interval? Exact Littlestone dimensions of intervals and their unions}.
Preprint \href{https://arxiv.org/abs/2611.28727}{arXiv:2611.28727}, November 2026.

\end{thebibliography}

\end{document}
